\documentclass[12pt, a4paper]{article}
% --- 1. Cấu hình Tiếng Việt và Font chữ chuẩn ---
\usepackage[utf8]{inputenc}
\usepackage[T5]{fontenc}
\usepackage[vietnamese]{babel}
\usepackage{mathptmx} % Font Times New Roman đồng bộ cả chữ và công thức toán, không gây xung đột gói

\makeatletter
\renewcommand{\normalsize}{\@setfontsize\normalsize{13pt}{16pt}}
\makeatother
\normalsize

% --- 2. Gói Toán học & Ký hiệu bổ trợ ---
\usepackage{amsmath, amssymb, amsfonts, mathrsfs, amsthm}
\usepackage{graphicx}
\usepackage{fontawesome}
\usepackage{booktabs, tabularx, array, multirow}
\usepackage{enumitem}

% --- 3. Đồ họa TikZ, Bảng biến thiên & Hình không gian ---
\usepackage{tikz}
\usetikzlibrary{calc, intersections, angles, quotes, patterns, arrows.meta, 3d}
\usepackage{pgfplots}
\pgfplotsset{compat=1.18}
\usepackage{tkz-euclide}
\usepackage{tkz-tab}

\renewcommand{\vec}[1]{\overrightarrow{#1}}

% --- 4. Hộp sư phạm (Tcolorbox) ---
\usepackage[many]{tcolorbox}
\tcbuselibrary{skins, breakable}

% --- 5. Căn lề chuẩn văn bản giáo án ---
\usepackage{geometry}
\geometry{
    a4paper,
    left=25mm,
    right=15mm,
    top=20mm,
    bottom=20mm
}

% --- 6. Header / Footer chuẩn hóa ---
\usepackage{fancyhdr}
\usepackage{lastpage}
\pagestyle{fancy}
\fancyhf{}

\fancyhead[L]{\small \textit{Kế hoạch bài dạy Toán 11}}
\fancyhead[R]{\textbf{Trang \thepage/\pageref{LastPage}}}
\fancyfoot[L]{\small \textit{Tổ Toán - Tin}}
\fancyfoot[R]{\small \textit{Trường THCS\&THPT Hồng Vân}}
\renewcommand{\headrulewidth}{0.4pt}
\renewcommand{\footrulewidth}{0.4pt}

% --- 7. Giãn dòng & Giãn đoạn theo chuẩn Nghị định 30/2020/NĐ-CP ---
\usepackage{setspace}
\setstretch{1.15}
\setlength{\parindent}{1.0cm}
\setlength{\parskip}{5pt}

% Định nghĩa hộp sư phạm chuyên dụng
\newtcolorbox{pedabox}[2][]{
    enhanced,
    breakable,
    colback=blue!3!white,
    colframe=blue!65!black,
    fonttitle=\bfseries\small,
    title={#2},
    arc=2mm,
    boxrule=0.6pt,
    left=3mm, right=3mm, top=2mm, bottom=2mm,
    #1
}

\newtcolorbox{aibox}[2][]{
    enhanced,
    breakable,
    colback=teal!4!white,
    colframe=teal!70!black,
    fonttitle=\bfseries\small,
    title={#2},
    arc=2mm,
    boxrule=0.6pt,
    left=3mm, right=3mm, top=2mm, bottom=2mm,
    #1
}

\newtcolorbox{scaffoldbox}[2][]{
    enhanced,
    breakable,
    colback=orange!4!white,
    colframe=orange!80!black,
    fonttitle=\bfseries\small,
    title={#2},
    arc=2mm,
    boxrule=0.6pt,
    left=3mm, right=3mm, top=2mm, bottom=2mm,
    #1
}

\newtcolorbox{stembox}[2][]{
    enhanced,
    breakable,
    colback=purple!4!white,
    colframe=purple!75!black,
    fonttitle=\bfseries\small,
    title={#2},
    arc=2mm,
    boxrule=0.6pt,
    left=3mm, right=3mm, top=2mm, bottom=2mm,
    #1
}

\begin{document}

\begin{center}
    {\fontsize{14pt}{17pt}\selectfont \textbf{KẾ HOẠCH BÀI DẠY MÔN TOÁN LỚP 11}}\\[6pt]
    {\fontsize{16pt}{19pt}\selectfont \textbf{BÀI TẬP CUỐI CHƯƠNG VI}}\\[4pt]
    {\fontsize{12pt}{15pt}\selectfont \textbf{Môn học: Toán 11} \ \textbar \ \textbf{Thời lượng: 1 tiết (Tiết 62 theo PPCT | Tuần 21)}}
\end{center}

\vspace{5pt}

\section*{I. MỤC TIÊU}

\subsection*{1. Kiến thức, Năng lực}

\begin{itemize}[leftmargin=1.2cm]
    \item \textbf{Yêu cầu cần đạt (Nguyên văn CT GDPT 2018 Toán 11):}
    \begin{itemize}[leftmargin=0.6cm]
        \item Hệ thống hoá, củng cố toàn bộ kiến thức trọng tâm của Chương VI: Phép tính lũy thừa với số mũ thực; Khái niệm và tính chất của lôgarit; Hàm số mũ và hàm số lôgarit; Phương pháp giải phương trình, bất phương trình mũ và lôgarit cơ bản.
        \item Củng cố và thành thạo phương pháp giải phương trình, bất phương trình mũ và lôgarit cơ bản (đặc biệt dạng đưa về cùng cơ số $a^{f(x)} = a^{g(x)}$ và $\log_a f(x) = \log_a g(x)$).
        \item Giải quyết được các bài toán thực tiễn liên quan đến hàm số mũ và hàm số lôgarit (bài toán tăng trưởng dân số, lãi kép ngân hàng, phân rã phóng xạ, độ pH hóa học, thang đo địa chấn Richter).
    \end{itemize}

    \item \textbf{Năng lực Toán học đặc thù:}
    \begin{itemize}[leftmargin=0.6cm]
        \item \textit{Năng lực tư duy và lập luận toán học:} Khái quát hoá mối liên hệ tương hỗ giữa lũy thừa và lôgarit; phân biệt bản chất biến thiên của hàm số mũ, lôgarit và quy tắc đổi chiều bất phương trình theo cơ số ($a > 1$ giữ nguyên chiều, $0 < a < 1$ đổi chiều); lập luận chặt chẽ về điều kiện xác định để loại bỏ nghiệm ngoại lai.
        \item \textit{Năng lực mô hình hóa toán học:} Thiết lập mô hình toán học sử dụng hàm số mũ hoặc lôgarit từ các tình huống kinh tế - xã hội, khoa học tự nhiên để giải bài toán tìm thời gian, số lượng hoặc tốc độ tăng trưởng.
        \item \textit{Năng lực giải quyết vấn đề toán học:} Lựa chọn chiến lược giải tối ưu: kết hợp các phép biến đổi đại số, quy tắc lôgarit hóa, mũ hóa, đặt ẩn phụ để giải quyết các bài toán tổng hợp chương.
        \item \textit{Năng lực giao tiếp toán học:} Trình bày hệ thống kiến thức bằng sơ đồ tư duy; sử dụng chính xác các ký hiệu toán học ($\iff, \implies, \in, \cap, \cup, \emptyset$), diễn đạt tập nghiệm mạch lạc.
        \item \textit{Năng lực sử dụng công cụ, phương tiện học toán:} Thành thạo máy tính cầm tay (MTCT) Casio fx-580VN X trong việc tính toán giá trị lũy thừa, lôgarit, dò tìm nghiệm phương trình bằng lệnh \texttt{SOLVE}, kiểm tra tập nghiệm bất phương trình bằng chức năng \texttt{CALC} và \texttt{TABLE}.
    \end{itemize}

    \item \textbf{Năng lực số (Theo Thông tư 02/2025/TT-BGDĐT):}
    \begin{itemize}[leftmargin=0.6cm]
        \item \textit{Mã 2.1NC1a (Tương tác và chia sẻ qua môi trường số):} Học sinh biết tương tác, chia sẻ các sản phẩm sơ đồ tư duy tổng kết chương và các dạng bài giải phương trình mũ, lôgarit qua nhóm học tập trực tuyến (Padlet, Google Classroom, nhóm Zalo học tập).
    \end{itemize}

    \item \textbf{Năng lực AI (Theo Quyết định 2422/QĐ-BGDĐT):}
    \begin{itemize}[leftmargin=0.6cm]
        \item \textit{Mã 11.B3.1 (Đạo đức, trách nhiệm và tính liêm chính trong học tập số với AI):} Học sinh thể hiện thái độ trung thực, liêm chính khi làm bài kiểm tra trắc nghiệm online; tuyệt đối không gian lận hay sao chép kết quả từ các công cụ AI giải toán trong quá trình làm bài; biết sử dụng AI đúng đắn sau khi hoàn thành bài thi như một trợ lý gia sư phân tích nguyên nhân sai sót của từng câu hỏi.
    \end{itemize}

    \item \textbf{Năng lực chung:}
    \begin{itemize}[leftmargin=0.6cm]
        \item \textit{Tự chủ và tự học:} Tự hệ thống hóa kiến thức đã học, tự đánh giá mức độ nắm vững từng đơn vị bài học trong chương.
        \item \textit{Giao tiếp và hợp tác:} Hoạt động nhóm hiệu quả trong việc trao đổi phương pháp giải và phản biện kết quả học tập.
    \end{itemize}
\end{itemize}

\subsection*{2. Phẩm chất}

\begin{itemize}[leftmargin=1.2cm]
    \item \textbf{Chăm chỉ:} Tích cực ôn luyện, hoàn thành đầy đủ hệ thống bài tập tổng kết chương.
    \item \textbf{Trung thực:} Tự giác, khách quan, liêm chính trong kiểm tra, đánh giá trực tuyến.
    \item \textbf{Trách nhiệm:} Có ý thức bảo vệ an toàn thông tin số trong nhóm học tập, sẵn sàng hỗ trợ bạn bè cùng tiến bộ.
\end{itemize}

\vspace{5pt}

\section*{II. THIẾT BỊ DẠY HỌC VÀ HỌC LIỆU}

\begin{itemize}[leftmargin=1.2cm]
    \item \textbf{Giáo viên:}
    \begin{itemize}[leftmargin=0.6cm]
        \item Kế hoạch bài dạy chi tiết theo chuẩn Công văn 5512/BGDĐT-GDTrH.
        \item Máy tính, máy chiếu, bảng tương tác, phần mềm giả lập MTCT Casio fx-580VN X.
        \item Hệ thống bài kiểm tra trắc nghiệm online (trên nền tảng Azota / Google Forms / Quizizz) đánh giá mức độ đạt chuẩn Chương VI.
        \item Phiếu học tập ôn tập tổng hợp phân hóa bậc thang (Phiếu số 1) và Bảng Rubric đánh giá năng lực học sinh.
    \end{itemize}
    \item \textbf{Học sinh:}
    \begin{itemize}[leftmargin=0.6cm]
        \item Sách giáo khoa Toán 11, vở bài tập, giấy nháp.
        \item Máy tính cầm tay Casio fx-580VN X (hoặc tương đương).
        \item Điện thoại thông minh hoặc máy tính cá nhân kết nối mạng Internet để tham gia làm bài kiểm tra trắc nghiệm online.
        \item Sơ đồ tư duy ôn tập Chương VI đã chuẩn bị trước ở nhà theo nhóm.
    \end{itemize}
\end{itemize}

\vspace{5pt}

\section*{III. TIẾN TRÌNH DẠY HỌC (1 TIẾT | TIẾT 62 THEO PPCT)}

\subsection*{\faClockO \ TIẾT 62 THEO PPCT: HỆ THỐNG HÓA KIẾN THỨC VÀ BÀI TẬP TỔNG KẾT CHƯƠNG VI}

\subsubsection*{1. Hoạt động 1: Mở đầu (Khởi động) (5 phút)}

\noindent \textbf{a) Mục tiêu:}
Tạo không khí học tập hào hứng, kích hoạt lại hệ thống định nghĩa và công thức cốt lõi của Chương VI thông qua trò chơi ghép cặp công thức nhanh.

\noindent \textbf{b) Nội dung:}
Giáo viên trình chiếu trò chơi tương tác \textbf{``Ghép đôi hoàn hảo - Lũy thừa \& Lôgarit''}:
\begin{pedabox}{Trò chơi khởi động: Ghép đôi công thức tương đương}
Nối mỗi biểu thức ở Cột A với biểu thức tương đương đúng ở Cột B ($a, b, c > 0; a, c \neq 1; \alpha, \beta \in \mathbb{R}$):
\begin{center}
\begin{tabular}{|l|l|}
\hline
\textbf{Cột A} & \textbf{Cột B} \\ \hline
1. $\log_a (b \cdot c)$ & A. $\alpha \cdot \log_a b$ \\ \hline
2. $\log_a \left(\dfrac{b}{c}\right)$ & B. $\dfrac{1}{\log_a b}$ ($b \neq 1$) \\ \hline
3. $\log_a (b^\alpha)$ & C. $\log_a b + \log_a c$ \\ \hline
4. $\log_b a$ & D. $\log_a b - \log_a c$ \\ \hline
5. $a^{\log_a b}$ & E. $b$ \\ \hline
\end{tabular}
\end{center}
\end{pedabox}

\noindent \textbf{c) Sản phẩm:}
Kết quả ghép nối chuẩn xác của học sinh:
\begin{center}
\textbf{1 - C} \quad \textbar \quad \textbf{2 - D} \quad \textbar \quad \textbf{3 - A} \quad \textbar \quad \textbf{4 - B} \quad \textbar \quad \textbf{5 - E}
\end{center}

\noindent \textbf{d) Tổ chức thực hiện:}
\begin{itemize}[leftmargin=0.6cm]
    \item \textbf{Giao nhiệm vụ:} Giáo viên trình chiếu màn hình ghép đôi, yêu cầu học sinh xung phong trả lời nhanh trong 2 phút.
    \item \textbf{Thực hiện nhiệm vụ:} 2 học sinh nhanh nhất trả lời miệng; các học sinh khác theo dõi, nhận xét.
    \item \textbf{Báo cáo, thảo luận:} Cả lớp đồng thanh kiểm tra lại các công thức quan trọng.
    \item \textbf{Kết luận, nhận định:} Giáo viên chốt đáp án, biểu dương tinh thần chuẩn bị bài và dẫn dắt vào tiết luyện tập tổng kết chương.
\end{itemize}

\subsubsection*{2. Hoạt động 2: Luyện tập - Hệ thống hóa kiến thức và Chữa bài tập trọng tâm (25 phút)}

\paragraph{2.1. Nhiệm vụ 1: Hệ thống hóa kiến thức Chương VI bằng Sơ đồ tư duy (8 phút)}

\noindent \textbf{a) Mục tiêu:}
Khái quát hóa toàn bộ kiến thức Chương VI dưới dạng sơ đồ logic; chia sẻ sản phẩm số trên nhóm học tập trực tuyến.

\noindent \textbf{b) Nội dung:}
Giáo viên cùng học sinh tổng kết 4 nhánh kiến thức lớn của Chương VI:
\begin{enumerate}[leftmargin=0.6cm]
    \item \textbf{Nhánh 1: Lũy thừa với số mũ thực:}
    \begin{itemize}
        \item Với $n \in \mathbb{Z}^+$: $a^n = \underbrace{a \cdot a \dots a}_{n\text{ thừa số}}$.
        \item Với $a \neq 0$: $a^0 = 1; a^{-n} = \dfrac{1}{a^n}$.
        \item Với $a > 0, m \in \mathbb{Z}, n \in \mathbb{N}, n \ge 2$: $a^{\frac{m}{n}} = \sqrt[n]{a^m}$.
        \item Các tính chất: $a^\alpha \cdot a^\beta = a^{\alpha + \beta}$; $\dfrac{a^\alpha}{a^\beta} = a^{\alpha - \beta}$; $(a^\alpha)^\beta = a^{\alpha \cdot \beta}$; $(a \cdot b)^\alpha = a^\alpha \cdot b^\alpha$.
    \end{itemize}
    \item \textbf{Nhánh 2: Lôgarit:}
    \begin{itemize}
        \item Định nghĩa: Với $a > 0, a \neq 1, b > 0$: $\alpha = \log_a b \iff a^\alpha = b$.
        \item Các công thức biến đổi cơ bản: Lôgarit của tích, thương, lũy thừa; công thức đổi cơ số $\log_a b = \dfrac{\log_c b}{\log_c a}$.
    \end{itemize}
    \item \textbf{Nhánh 3: Hàm số mũ $y = a^x$ và Hàm số lôgarit $y = \log_a x$ ($a > 0, a \neq 1$):}
    \begin{itemize}
        \item Tập xác định, tập giá trị: Hàm số mũ có $D = \mathbb{R}, T = (0; +\infty)$; Hàm số lôgarit có $D = (0; +\infty), T = \mathbb{R}$.
        \item Chiều biến thiên: Đồng biến trên tập xác định khi $a > 1$; Nghịch biến trên tập xác định khi $0 < a < 1$.
        \item Đồ thị: Nhận trục $Ox$ làm tiệm cận ngang (hàm mũ); nhận trục $Oy$ làm tiệm cận đứng (hàm lôgarit).
    \end{itemize}
    \item \textbf{Nhánh 4: Phương trình và Bất phương trình mũ, lôgarit:}
    \begin{itemize}
        \item Phương pháp đưa về cùng cơ số: $a^{f(x)} = a^{g(x)} \iff f(x) = g(x)$; $\log_a f(x) = \log_a g(x) \iff \begin{cases} f(x) > 0 \\ f(x) = g(x) \end{cases}$.
        \item Bất phương trình: $a > 1 \implies$ giữ nguyên chiều; $0 < a < 1 \implies$ bắt buộc đổi chiều.
    \end{itemize}
\end{enumerate}

\begin{aibox}{Tích hợp Năng lực số (Theo Thông tư 02/2025/TT-BGDĐT - Mã 2.1NC1a)}
\textbf{Hoạt động tương tác và chia sẻ sơ đồ tư duy trên môi trường số:}
\begin{itemize}[leftmargin=0.6cm]
    \item Đại diện 4 nhóm học sinh tải sản phẩm sơ đồ tư duy (Mindmap số vẽ bằng Canva / GitMind / XMind) lên bức tường số Padlet của lớp theo đường link giáo viên cung cấp.
    \item Học sinh sử dụng thiết bị số để bình chọn (thả tim), bình luận nhận xét tính đầy đủ, sáng tạo và chuẩn xác về công thức toán học giữa các nhóm.
\end{itemize}
\end{aibox}

\begin{scaffoldbox}{Giàn giáo sư phạm: Bảng so sánh phương trình và bất phương trình mũ vs lôgarit}
\begin{center}
\begin{tabular}{|m{2.5cm}|m{5.5cm}|m{6cm}|}
\hline
\textbf{Đặc trưng} & \textbf{Dạng Mũ ($a > 0, a \neq 1$)} & \textbf{Dạng Lôgarit ($a > 0, a \neq 1$)} \\ \hline
\textbf{Điều kiện xác định} & Số mũ $f(x)$ xác định theo tính chất của biểu thức $f(x)$. & \textbf{BẮT BUỘC} biểu thức trong dấu lôgarit phải dương: $f(x) > 0$. \\ \hline
\textbf{Đưa về cùng cơ số} & $a^{f(x)} = a^{g(x)} \iff f(x) = g(x)$ & $\log_a f(x) = \log_a g(x) \iff \begin{cases} f(x) > 0 \\ f(x) = g(x) \end{cases}$ \\ \hline
\textbf{Bất phương trình ($a > 1$)} & $a^{f(x)} > a^{g(x)} \iff f(x) > g(x)$ (Giữ chiều) & $\log_a f(x) > \log_a g(x) \iff f(x) > g(x) > 0$ (Giữ chiều) \\ \hline
\textbf{Bất phương trình ($0 < a < 1$)} & $a^{f(x)} > a^{g(x)} \iff f(x) < g(x)$ (\textbf{Đổi chiều}) & $\log_a f(x) > \log_a g(x) \iff 0 < f(x) < g(x)$ (\textbf{Đổi chiều}) \\ \hline
\end{tabular}
\end{center}
\end{scaffoldbox}

\paragraph{2.2. Nhiệm vụ 2: Luyện giải các dạng bài tập trọng tâm Chương VI (17 phút)}

\noindent \textbf{a) Mục tiêu:}
Rèn luyện kỹ năng giải các bài toán tổng hợp: Rút gọn biểu thức, tìm tập xác định, giải phương trình và bất phương trình mũ, lôgarit cơ bản.

\noindent \textbf{b) Nội dung:}
Giáo viên giao phiếu bài tập:
\begin{pedabox}{Bài tập rèn luyện kỹ năng trọng tâm}
\begin{enumerate}[leftmargin=0.6cm]
    \item \textbf{Dạng 1 (Biến đổi lũy thừa và lôgarit):} Rút gọn biểu thức:
    $$A = \dfrac{a^{\frac{1}{3}} \sqrt{b} + b^{\frac{1}{3}} \sqrt{a}}{\sqrt[6]{a} + \sqrt[6]{b}} \quad (a, b > 0); \qquad B = \log_2 3 \cdot \log_3 4 \cdot \log_4 5 \dots \log_{31} 32$$
    \item \textbf{Dạng 2 (Phương trình mũ và lôgarit):} Giải các phương trình:
    $$a) \ 4^{x^2 - 2x} = 2^{3x - 2}; \qquad b) \ \log_2 (x + 1) + \log_2 (x - 1) = 3$$
    \item \textbf{Dạng 3 (Bất phương trình mũ và lôgarit):} Giải các bất phương trình:
    $$a) \ \left(\dfrac{1}{3}\right)^{2x - 1} \ge 9; \qquad b) \ \log_{0{,}5} (2x - 3) > -2$$
\end{enumerate}
\end{pedabox}

\noindent \textbf{c) Sản phẩm:}
Lời giải chi tiết từng bài toán:
\begin{itemize}[leftmargin=0.6cm]
    \item \textbf{Dạng 1:}
    \begin{itemize}
        \item Ta có: $a^{\frac{1}{3}} = a^{\frac{2}{6}} = \left(\sqrt[6]{a}\right)^2$; $\sqrt{b} = b^{\frac{3}{6}} = \left(\sqrt[6]{b}\right)^3$.
        Tử số: $a^{\frac{1}{3}} \sqrt{b} + b^{\frac{1}{3}} \sqrt{a} = \left(\sqrt[6]{a}\right)^2 \left(\sqrt[6]{b}\right)^3 + \left(\sqrt[6]{b}\right)^2 \left(\sqrt[6]{a}\right)^3 = \left(\sqrt[6]{a}\right)^2 \left(\sqrt[6]{b}\right)^2 (\sqrt[6]{b} + \sqrt[6]{a})$.
        Vậy $A = \dfrac{(\sqrt[6]{a} \cdot \sqrt[6]{b})^2 (\sqrt[6]{a} + \sqrt[6]{b})}{\sqrt[6]{a} + \sqrt[6]{b}} = (\sqrt[6]{ab})^2 = (ab)^{\frac{2}{6}} = (ab)^{\frac{1}{3}} = \sqrt[3]{ab}$.
        \item Áp dụng công thức đổi cơ số liên tiếp:
        $$B = \log_2 3 \cdot \log_3 4 \dots \log_{31} 32 = \log_2 32 = \log_2 (2^5) = 5.$$
    \end{itemize}

    \item \textbf{Dạng 2:}
    \begin{itemize}
        \item $a)$ Ta có $4^{x^2 - 2x} = (2^2)^{x^2 - 2x} = 2^{2x^2 - 4x}$.
        Phương trình $\iff 2^{2x^2 - 4x} = 2^{3x - 2} \iff 2x^2 - 4x = 3x - 2 \iff 2x^2 - 7x + 2 = 0$.
        Giải ra ta được: $x_1 = \dfrac{7 + \sqrt{33}}{4}$; $x_2 = \dfrac{7 - \sqrt{33}}{4}$.
        Tập nghiệm là $S = \left\{\dfrac{7 - \sqrt{33}}{4}; \dfrac{7 + \sqrt{33}}{4}\right\}$.
        \item $b)$ Điều kiện: $\begin{cases} x + 1 > 0 \\ x - 1 > 0 \end{cases} \iff x > 1$.
        Phương trình $\iff \log_2 [(x + 1)(x - 1)] = 3 \iff x^2 - 1 = 2^3 = 8 \iff x^2 = 9 \iff \left[\begin{array}{l} x = 3 \\ x = -3 \end{array}\right.$.
        Đối chiếu điều kiện $x > 1$, ta loại $x = -3$, nhận $x = 3$. Tập nghiệm là $S = \{3\}$.
    \end{itemize}

    \item \textbf{Dạng 3:}
    \begin{itemize}
        \item $a)$ Vì $9 = 3^2 = \left(\dfrac{1}{3}\right)^{-2}$ và cơ số $\dfrac{1}{3} \in (0; 1)$ nên \textbf{đổi chiều}:
        $$\left(\dfrac{1}{3}\right)^{2x - 1} \ge \left(\dfrac{1}{3}\right)^{-2} \iff 2x - 1 \le -2 \iff 2x \le -1 \iff x \le -\dfrac{1}{2}.$$
        Tập nghiệm là $S = \left(-\infty; -\dfrac{1}{2}\right]$.
        \item $b)$ Điều kiện: $2x - 3 > 0 \iff x > \dfrac{3}{2}$.
        Vì cơ số $0{,}5 \in (0; 1)$ nên \textbf{đổi chiều}:
        $$\log_{0{,}5} (2x - 3) > -2 \iff 2x - 3 < (0{,}5)^{-2} = \left(\dfrac{1}{2}\right)^{-2} = 4 \iff 2x < 7 \iff x < \dfrac{7}{2}.$$
        Kết hợp điều kiện $x > \dfrac{3}{2}$, ta được: $\dfrac{3}{2} < x < \dfrac{7}{2}$.
        Tập nghiệm là $S = \left(\dfrac{3}{2}; \dfrac{7}{2}\right)$.
    \end{itemize}
\end{itemize}

\noindent \textbf{d) Tổ chức thực hiện:}
\begin{itemize}[leftmargin=0.6cm]
    \item \textbf{Giao nhiệm vụ:} Chia lớp thành 3 nhóm tương ứng với 3 dạng bài. Học sinh làm bài độc lập trong 8 phút, sau đó thảo luận cặp đôi.
    \item \textbf{Thực hiện nhiệm vụ:} Giáo viên theo dõi, hướng dẫn học sinh yếu cách đưa về cơ số 2 ở Dạng 2a, và nhắc nhở đổi chiều bất phương trình ở Dạng 3.
    \item \textbf{Báo cáo, thảo luận:} 3 học sinh đại diện lên bảng trình bày 3 dạng; cả lớp nhận xét, phản biện.
    \item \textbf{Kết luận, nhận định:} Giáo viên chuẩn hóa lời giải và lưu ý các bẫy thường gặp trong đề thi tốt nghiệp THPT.
\end{itemize}

\subsubsection*{3. Hoạt động 3: Vận dụng thực tiễn và Kiểm tra trực tuyến (10 phút)}

\noindent \textbf{a) Mục tiêu:}
Vận dụng hàm số mũ giải quyết bài toán mô hình tăng trưởng dân số thực tế; thực hiện bài kiểm tra trắc nghiệm online rèn luyện tính trung thực và liêm chính học thuật số.

\noindent \textbf{b) Nội dung:}
\begin{stembox}{Bài toán Mô hình hóa Tăng trưởng Dân số (STEM)}
Dân số của một tỉnh vào đầu năm 2024 là $P_0 = 1{,}5$ triệu người. Giả sử tỷ lệ tăng dân số hàng năm của tỉnh đó ổn định ở mức $r = 1{,}1\%$/năm. Theo mô hình tăng trưởng hàm số mũ liên tục:
$$P(t) = P_0 \cdot e^{r \cdot t}$$
(trong đó $t$ là số năm tính từ đầu năm 2024).
\begin{enumerate}[leftmargin=0.6cm]
    \item Dự báo dân số của tỉnh đó vào đầu năm 2030 ($t = 6$).
    \item Hỏi bắt đầu từ năm nào thì dân số của tỉnh đó vượt quá $2{,}0$ triệu người?
\end{enumerate}
\end{stembox}

\begin{aibox}{Tích hợp Năng lực AI (Theo QĐ 2422/QĐ-BGDĐT - Mã 11.B3.1)}
\textbf{Thực hiện kiểm tra trắc nghiệm số và Liêm chính học thuật:}
\begin{itemize}[leftmargin=0.6cm]
    \item \textbf{Nhiệm vụ kiểm tra:} Học sinh quét mã QR trên màn hình tham gia bài kiểm tra trắc nghiệm ngắn 5 câu (thời gian 4 phút) trên hệ thống Azota/Google Forms.
    \item \textbf{Cam kết liêm chính:} Học sinh làm bài hoàn toàn tự lực, tuyệt đối không sử dụng công cụ AI để sao chép đáp án.
    \item \textbf{Ứng dụng AI sau khi nộp bài:} Giáo viên hướng dẫn học sinh xuất kết quả làm bài, nhập prompt vào trợ lý AI cá nhân: \textit{``Giải thích chi tiết vì sao câu 4 em chọn đáp án B là sai và hướng dẫn phương pháp giải đúng bằng kiến thức đổi chiều bất phương trình lôgarit cơ số nhỏ hơn 1.''}
\end{itemize}
\end{aibox}

\noindent \textbf{c) Sản phẩm:}
Lời giải bài toán dân số:
\begin{enumerate}[leftmargin=0.6cm]
    \item Vào đầu năm 2030, tức là sau $t = 2030 - 2024 = 6$ năm:
    $$P(6) = 1{,}5 \cdot e^{0{,}011 \cdot 6} = 1{,}5 \cdot e^{0{,}066} \approx 1{,}6023\text{ (triệu người)} = 1\,602\,300\text{ người}.$$
    \item Để dân số vượt quá $2{,}0$ triệu người:
    $$P(t) > 2{,}0 \iff 1{,}5 \cdot e^{0{,}011 \cdot t} > 2{,}0 \iff e^{0{,}011 \cdot t} > \dfrac{2{,}0}{1{,}5} = \dfrac{4}{3}$$
    Lấy lôgarit tự nhiên ($\ln$) hai vế (vì cơ số $e \approx 2{,}718 > 1$ giữ nguyên chiều):
    $$0{,}011 \cdot t > \ln\left(\dfrac{4}{3}\right) \iff t > \dfrac{\ln(4/3)}{0{,}011} \approx \dfrac{0{,}28768}{0{,}011} \approx 26{,}15\text{ năm}.$$
    Như vậy, sau hơn $26$ năm, tức là bắt đầu từ năm $2024 + 27 = 2051$, dân số của tỉnh đó sẽ vượt quá $2$ triệu người.
\end{enumerate}

\noindent \textbf{d) Tổ chức thực hiện:}
\begin{itemize}[leftmargin=0.6cm]
    \item \textbf{Giao nhiệm vụ:} Giáo viên giao bài toán STEM dân số và hướng dẫn học sinh bấm phím $\ln$ và $e^x$ trên MTCT.
    \item \textbf{Thực hiện nhiệm vụ:} Học sinh tính toán, thảo luận kết quả làm tròn năm nguyên dương. Sau đó, tiến hành làm bài kiểm tra trắc nghiệm online 4 phút.
    \item \textbf{Báo cáo, nhận định:} Hệ thống tự động chấm điểm và công bố phổ điểm; giáo viên tuyên dương các học sinh đạt điểm cao và trung thực trong kiểm tra.
\end{itemize}

\subsubsection*{4. Hoạt động 4: Tìm tòi, mở rộng và Hướng dẫn tự học (5 phút)}

\noindent \textbf{a) Mục tiêu:}
Định hướng học sinh tự ôn tập củng cố toàn bộ Chương VI và chuẩn bị bài mới cho Chương VII: Quan hệ vuông góc trong không gian.

\noindent \textbf{b) Nội dung:}
\begin{itemize}[leftmargin=0.6cm]
    \item Tự hoàn thiện sơ đồ tư duy cá nhân toàn bộ Chương VI vào vở ghi.
    \item Ôn tập lại định nghĩa vectơ trong không gian và giá của vectơ (đã học ở lớp 10) để chuẩn bị cho \textbf{Bài 22. Hai đường thẳng vuông góc} (Chương VII).
    \item Tìm hiểu trước cách xác định góc giữa hai đường thẳng trong không gian bằng phần mềm GeoGebra 3D.
\end{itemize}

\noindent \textbf{c) Sản phẩm:}
Kế hoạch học tập cá nhân và bài tập về nhà trong vở bài tập của học sinh.

\noindent \textbf{d) Tổ chức thực hiện:}
Giáo viên hướng dẫn, giao nhiệm vụ và căn dặn cả lớp về nhà chuẩn bị đồ dùng học hình học không gian (thước kẻ, bút màu, compa).

\vspace{10pt}

\section*{IV. PHỤ LỤC HỌC LIỆU BỔ TRỢ (BẮT BUỘC)}

\subsection*{1. Phiếu học tập số 1: Ôn tập tổng hợp Chương VI}

\begin{pedabox}{PHIẾU HỌC TẬP TỔNG HỢP CHƯƠNG VI}
\textbf{Họ và tên học sinh:} \dotfill \ \textbf{Lớp:} \dotfill \ \textbf{Điểm:} \dotfill

\noindent \textbf{PHẦN 1: TRẮC NGHIỆM NHANH (4 ĐIỂM)}
\begin{enumerate}[leftmargin=0.6cm]
    \item Với $a > 0$ và $a \neq 1$, giá trị của biểu thức $\log_a (a^3)$ bằng:
    \begin{center}
    \textbf{A.} $1$ \qquad \qquad \textbf{B.} $3$ \qquad \qquad \textbf{C.} $\dfrac{1}{3}$ \qquad \qquad \textbf{D.} $a^3$
    \end{center}
    \item Tập xác định của hàm số $y = \log_2 (3x - 6)$ là:
    \begin{center}
    \textbf{A.} $D = [2; +\infty)$ \qquad \textbf{B.} $D = (2; +\infty)$ \qquad \textbf{C.} $D = (-\infty; 2)$ \qquad \textbf{D.} $D = \mathbb{R} \setminus \{2\}$
    \end{center}
    \item Hàm số nào sau đây đồng biến trên $\mathbb{R}$?
    \begin{center}
    \textbf{A.} $y = \left(\dfrac{1}{2}\right)^x$ \qquad \textbf{B.} $y = (0{,}8)^x$ \qquad \textbf{C.} $y = (\sqrt{3})^x$ \qquad \textbf{D.} $y = \left(\dfrac{\pi}{4}\right)^x$
    \end{center}
    \item Nghiệm của phương trình $3^{x+1} = 9$ là:
    \begin{center}
    \textbf{A.} $x = 1$ \qquad \qquad \textbf{B.} $x = 2$ \qquad \qquad \textbf{C.} $x = 3$ \qquad \qquad \textbf{D.} $x = 0$
    \end{center}
\end{enumerate}

\noindent \textbf{PHẦN 2: TỰ LUẬN TỔNG HỢP (6 ĐIỂM)}
\begin{enumerate}[leftmargin=0.6cm]
    \item \textbf{Bài 1 (2 điểm):} Cho $\log_2 3 = a$ và $\log_2 5 = b$. Hãy biểu diễn $\log_2 45$ theo $a$ và $b$.
    \item \textbf{Bài 2 (2 điểm):} Giải bất phương trình: $\log_{\frac{1}{2}} (x^2 - 5x + 7) \ge 0$.
    \item \textbf{Bài 3 (2 điểm):} Độ pH của một nguồn nước sinh hoạt được xác định bởi công thức $\text{pH} = -\log[\text{H}^+]$. Theo quy chuẩn quốc gia, nước sinh hoạt an toàn khi độ pH nằm trong khoảng từ $6{,}5$ đến $8{,}5$. Hãy xác định khoảng an toàn của nồng độ mol ion $[\text{H}^+]$.
\end{enumerate}
\end{pedabox}

\subsection*{2. Bảng Rubric đánh giá năng lực học sinh trong tiết học}

\begin{center}
\begin{tabularx}{\textwidth}{|p{2.8cm}|X|X|X|X|}
    \hline
    \textbf{Tiêu chí} & \textbf{Mức 1: Chưa đạt} & \textbf{Mức 2: Đạt} & \textbf{Mức 3: Khá} & \textbf{Mức 4: Tốt} \\ \hline
    \textbf{1. Hệ thống hóa kiến thức chương} & Nhớ rời rạc một số công thức, hay nhầm dấu khi tính lũy thừa, lôgarit. & Tái hiện được các công thức cơ bản; vẽ được sơ đồ tư duy đơn giản. & Vẽ sơ đồ tư duy logic, đầy đủ 4 nhánh kiến thức; nắm vững công thức biến đổi. & Sơ đồ tư duy sáng tạo, chuyên nghiệp; liên kết sâu sắc giữa hàm mũ và lôgarit. \\ \hline
    \textbf{2. Kỹ năng giải bài tập tổng hợp} & Thường xuyên quên điều kiện xác định; quên đổi chiều khi cơ số nhỏ hơn 1. & Giải được bài tập ở mức độ cơ bản; biết đặt điều kiện biểu thức dương. & Giải thành thạo bài tập đưa về cùng cơ số; kết hợp nghiệm chính xác trên trục số. & Biến đổi linh hoạt bài toán phân hóa; giải nhanh và chính xác các dạng bài thực tiễn. \\ \hline
    \textbf{3. Năng lực số \& Tương tác số (Mã 2.1NC1a)} & Chưa biết cách chia sẻ sản phẩm số lên Padlet/LMS lớp học. & Tải được tệp ảnh sơ đồ tư duy lên nhóm học tập trực tuyến. & Thao tác thuần thục; tương tác tích cực, đóng góp ý kiến nhận xét nhóm bạn. & Quản trị tốt nội dung chia sẻ; sử dụng phần mềm số trình bày bài toán đẹp mắt. \\ \hline
    \textbf{4. Đạo đức số \& Liêm chính AI (Mã 11.B3.1)} & Còn phụ thuộc vào ứng dụng giải bài hoặc chưa trung thực trong bài test online. & Hoàn thành bài kiểm tra online tự lực; không sao chép đáp án từ AI. & Thể hiện thái độ trung thực tuyệt đối; biết dùng AI hỏi giải thích các câu sai. & Có tư duy phản biện cao; làm chủ AI có trách nhiệm; lan tỏa tinh thần liêm chính học thuật. \\ \hline
\end{tabularx}
\end{center}

\end{document}
